% Generated draft: requires mathematical and source review.
\documentclass[11pt]{article}
\usepackage[margin=25mm]{geometry}
\usepackage{fontspec}
\setmainfont{Latin Modern Roman}
\usepackage{amsmath,amssymb,mathtools,unicode-math}
\setmathfont{Latin Modern Math}
\usepackage{enumitem,needspace,xurl,hyperref}
\setlength{\parindent}{0pt}
\setlength{\parskip}{5pt}
\setlength{\emergencystretch}{3em}
\newcommand{\N}{\mathbb{N}}
\newcommand{\Z}{\mathbb{Z}}
\newcommand{\Q}{\mathbb{Q}}
\newcommand{\R}{\mathbb{R}}
\newcommand{\C}{\mathbb{C}}
\newcommand{\F}{\mathbb{F}}
\newcommand{\A}{\mathbb{A}}
\newcommand{\PP}{\mathbb{P}}
\newcommand{\E}{\mathbb{E}}
\newcommand{\eps}{\varepsilon}
\newcommand{\dd}{\,\mathrm{d}}
\begin{document}
\section*{UnsolvedMath ID: 30002005}
\textbf{Source problem code:} catalogue entry
\subsection{Definitions and mathematical statement}
The supplied background identifies the intended scope as all wonderful varieties, but does not specify their ground field, characteristic, acting group, or precise class. Let \(X\) range over that intended class. Write \(\Sigma_X\) for the unspecified set denoted by \(\Sigma\) in the source. Its formal additive monoid is
\[
\mathbb N[\Sigma_X]=\left\{\sum_{\sigma\in\Sigma_X}n_\sigma\sigma:\ n_\sigma\in\mathbb Z_{\geq 0},\quad n_\sigma=0\text{ for all but finitely many }\sigma\right\}.
\]
Here the use of nonnegative integer coefficients is an explicit convention for the formal monoid notation. Let \(\mathcal C_X\subseteq\mathbb N[\Sigma_X]\) denote the elements called covering relations in the source, and retain its notation \(\operatorname{ht}^{+}\) for positive height. These are placeholders: the supplied record does not define the covering relations, their underlying ordered set, or the positive-height function.

The requested bound, with the universal scope indicated by the supplied background, is
\[
\forall X\ \forall\gamma\in\mathcal C_X,\qquad \operatorname{ht}^{+}(\gamma)\leq 2.
\]
This is a statement schema requiring completion from the original mathematical context. In particular, it cannot be interpreted as a fully specified conjecture until the class of \(X\), the meaning of \(\Sigma_X\), the covering relation, and the definition and normalization of positive height are supplied. No identification of positive height with the coefficient sum in \(\mathbb N[\Sigma_X]\) is justified by this record.
\subsection{Short English statement}
The proposed universal bound says that every covering relation associated with a wonderful variety has positive height at most \(2\). The supplied record leaves the relevant structures and height function undefined, so the precise conjecture requires clarification from its source context.
\subsection{Sources}
\begin{itemize}
\item Local UnsolvedMath registry: supplied complete record for permanent ID \(30002005\), title ``Height Bounds for Covering Relations in Wonderful Varieties,'' exact source problem code \texttt{catalogue entry}, in the Oberwolfach Reports Open Problems collection. The supplied proposed year is 2012; no individual proposer is identified.
\item Supplied source citation: \emph{Enveloping Algebras and Geometric Representation Theory} (2012), Oberwolfach Reports, pp.\ 733--809, DOI 10.4171/owr/2012/13. The literature-source label is ``Oberwolfach Report 13 2012.'' \url{https://doi.org/10.4171/owr/2012/13}.
\item Supplied additional literature source: \emph{Projective Normality of Model Wonderful Varieties}. \url{https://doi.org/10.1090/ert/477}.
\item Literature assessment checked at \texttt{2026-08-21}: this is supplied dated curation, not independently verified. It labels the problem open and reports that projective-normality multiplication and the associated height condition are proved for model wonderful varieties and other special cases, while the universal second-condition conjecture for all wonderful varieties remains unresolved. The supplied background describes this assessment as literature triage derived from OWR curation and requests independent verification of time-sensitive status and cited claims.
\item The registry labels the statement \texttt{corrected\_verified} and reports a manual semantic, typographic, and LaTeX review claiming that the cleanup is self-contained and adds no unsupported source-dependent content. These are supplied review assertions; they do not resolve the missing mathematical definitions in the supplied record.
\end{itemize}
\end{document}
